\subsection*{Additional cluster and carbonyl data}

\begin{table}[H]
\begin{center}
\caption{Electronic energies, zero-point vibrational energies (ZPVE), and
relative energies for the \ce{Pt10} structures optimized
directly from the \ce{Pt5Ge5} parent framework. Each entry corresponds to its
own relaxed geometry and frequency calculation at the PBEPBE/def2-TZVP level.
$\Delta E_0$ is obtained from $E_0=E_{\rm elec}+\mathrm{ZPVE}$, taking the
orbital-stable triplet structure as the zero of the relative-energy scale.
}
\label{tab:tfg_pt10_spin_energies}
%\resizebox{\linewidth}{!}{%
\begin{tabular}{lrrrr}
\toprule
{$2S+1$} &
{$E_{\rm elec}$ / \(\mathrm{E_h}\)} & {ZPVE / \(\mathrm{E_h}\)} &
{$E_0$ / \(\mathrm{E_h}\)} & {$\Delta E_0$ / eV} \\
\midrule
3 & -1194.29464704 & 0.006301 & -1194.288346 & 0.0000 \\
5 & -1194.29312331 & 0.006296 & -1194.286827 & 0.0413 \\
7 & -1194.28876904 & 0.006134 & -1194.282635 & 0.1554 \\
9 & -1194.28126572 & 0.006102 & -1194.275164 & 0.3587 \\
\bottomrule
\end{tabular}
%}
\end{center}
\end{table}

\begin{table}[H]
\begin{center}
\caption{Summary of the AdNDP decompositions for the crossed-parent models.
Each model occupies one table row; multiline cells give the AdNDP elements,
occupation range, and interpretation. ON denotes the AdNDP occupation number.
For the triplet \ce{Pt10}, the alpha and beta density matrices were analyzed
independently. The 3c elements involving the same atomic centers in both spin
channels were paired to obtain their combined occupations.}
\label{tab:tfg_adndp_summary}
\renewcommand{\arraystretch}{1.08}
\setlength{\tabcolsep}{2.8pt}
\small
\begin{tabular}{@{}>{\raggedright\arraybackslash}p{1.25cm}
                 >{\raggedright\arraybackslash}p{1.55cm}
                 >{\raggedright\arraybackslash}p{5.25cm}
                 >{\raggedright\arraybackslash}p{3.75cm}
                 >{\raggedright\arraybackslash}p{2.75cm}@{}}
\toprule
 & {$2S+1$} & {AdNDP elements} & {ON range} & {Interpretation} \\
\midrule
			\ce{Pt10} & 3 & \begin{tabular}[t]{@{}l@{}}
			\(\alpha\): 30 $\times$ Pt-centered 1c\\
			\hspace{.3cm}21 $\times$ Pt 3c\\
			\(\beta\): 30 $\times$ Pt-centered 1c\\
			\hspace{.3cm}19 $\times$ Pt 3c
			\end{tabular} & \begin{tabular}[t]{@{}l@{}}
			\(\alpha\): 0.967--0.992 (1c)\\
			\hspace{.3cm}0.990--0.997 (3c)\\
			\(\beta\): 0.961--0.989 (1c)\\
			\hspace{.3cm}0.994--0.997 (3c)\\
			paired 3c: 1.984--1.994
			\end{tabular} & \begin{tabular}[t]{@{}l@{}}19 $\times$ 3c--2e\\
		2 $\times$ 3c--1e;\\delocalized\\multicenter model\end{tabular} \\\hline
\ce{Pt5Ge5} & 1 & \begin{tabular}[t]{@{}l@{}}10 $\times$ Pt 1c--2e\\2 $\times$ Pt--Pt 2c--2e\\22 $\times$ Pt--Ge 2c--2e\\1 $\times$ 4c--2e\end{tabular} & 1.89--1.99 & mainly covalent with residual multicenter bonding \\
\ce{Pt6Ge4} & 1 & \begin{tabular}[t]{@{}l@{}}12 $\times$ Pt 1c--2e\\5 $\times$
Pt--Pt 2c--2e\\21 $\times$ Pt--Ge 2c--2e\end{tabular} & 1.89--1.98 & localized
Pt--Pt/Pt--Ge bonding \\\hline
\ce{Pt4Ge6} & 1 & \begin{tabular}[t]{@{}l@{}}12 $\times$ Pt 1c--2e\\2 $\times$
Pt--Pt 2c--2e\\15 $\times$ Pt--Ge 2c--2e\\3 $\times$ Ge--Ge 2c--2e\end{tabular}
			& 1.81--1.98 & localized two-center bonding \\
			\bottomrule
\end{tabular}
\parbox{0.96\linewidth}{\footnotesize For \ce{Pt10}, the listed occupations
refer to individual spin channels and are therefore close to one. The residual
alpha and beta densities are 0.7131 and 0.8732 electrons, respectively. The
threshold parameter was 1.040 for the 1c and 3c searches, corresponding to a
minimum accepted spin-resolved ON of 0.960. Nineteen alpha and beta 3c
elements involve identical sets of Pt atoms; their summed occupations are
listed as paired 3c values. The remaining alpha 3c elements are located on
Pt(2)--Pt(3)--Pt(7) and Pt(5)--Pt(6)--Pt(10), with ON values of 0.99445 and
0.99450, respectively. For \ce{Pt5Ge5}, the threshold
parameters were 0.113, 0.135, and 0.200 for the 1c, 2c, and 4c searches,
respectively. The 1c and 2c parameters were 0.100 and 0.140 for \ce{Pt6Ge4},
and 0.160 and 0.250 for \ce{Pt4Ge6}.}
\end{center}
\end{table}

\begin{table}[H]
\begin{center}
\caption{Binding energies for \ce{CO} placed at different Pt atoms, together
with structural and vibrational descriptors, for carbonyls initiated from the
reported global-minimum (GM) parent structures. Each composition is represented
using the spin state of its bare-cluster GM. \protect\blue{The GM is a triplet for
\ce{Pt6Ge4} and a singlet for \ce{Pt5Ge5} and \ce{Pt4Ge6}.} Each $BE$ uses free singlet \ce{CO} and the
optimized bare GM cluster with the same spin state as reference.}
\label{tab:tfg_gm_bimetallic_adsorption_descriptors}
%\scriptsize
\begin{tabular}{lcrrrr}
\toprule
{Composition} & {Site} & {$BE$ / eV} &
{$r(\ce{Pt-C})$ / \AA} & {$r(\ce{C-O})$ / \AA} &
{$\nu(\ce{CO})$ / \si{\per\centi\metre}} \\
\midrule
\ce{Pt6Ge4CO} & 2 & -1.947 & 1.851 & 1.157 & 2018.5 \\
\ce{Pt6Ge4CO} & 6 & -1.842 & 1.853 & 1.157 & 2022.5 \\
\ce{Pt6Ge4CO} & 3 & -1.732 & 1.859 & 1.157 & 2016.6 \\
\ce{Pt6Ge4CO} & 1 & -1.657 & 1.852 & 1.157 & 2021.5 \\
\ce{Pt6Ge4CO} & 5 & -1.215 & 1.926 & 1.154 & 2011.2 \\
\ce{Pt6Ge4CO} & 4 & -1.207 & 1.909 & 1.155 & 2013.6 \\
\midrule
\ce{Pt5Ge5CO} & 1/2/4/5 & -1.346 & 1.898 & 1.157 & 2006.1 \\
\ce{Pt5Ge5CO} & 3 & -0.788 & 1.956 & 1.151 & 2022.0 \\
\midrule
\ce{Pt4Ge6CO} & 2 & -1.506 & 1.910 & 1.155 & 2009.6 \\
\ce{Pt4Ge6CO} & 4 & -1.240 & 1.927 & 1.154 & 2012.0 \\
\ce{Pt4Ge6CO} & 1 & -1.150 & 1.917 & 1.155 & 2006.7 \\
\ce{Pt4Ge6CO} & 3 & -1.046 & 1.915 & 1.155 & 2002.7 \\
\bottomrule
\end{tabular}
\end{center}
\end{table}

\begin{table}[H]
\begin{center}
\caption{Pre-adsorption local Hoppe ECoN descriptors for the Pt sites in the
optimized global-minimum (GM) Pt--Ge clusters. These values were calculated
to examine whether the local-environment patterns found for the
crossed-parent structures persist for the GM parents. Only metal neighbors
are included, with
$ECoN_{\rm total}^{0}=ECoN_{\rm Pt}^{0}+ECoN_{\rm Ge}^{0}$ and
$x_{\rm Ge}^{0}=ECoN_{\rm Ge}^{0}/ECoN_{\rm total}^{0}$. The binding
energies correspond to carbonyls with the spin state of the corresponding
bare-cluster GM. \protect{The GM is a triplet for \ce{Pt6Ge4} and a singlet for
\ce{Pt5Ge5} and \ce{Pt4Ge6}.}}
\label{tab:tfg_gm_ptge_econ_pre_descriptors}
\small
\begin{tabular}{lccrrrr}
\toprule
{System} & {Site} & {$BE$ / eV} &
{$ECoN_{\rm total}^{0}$} & {$ECoN_{\rm Pt}^{0}$} &
{$ECoN_{\rm Ge}^{0}$} & {$x_{\rm Ge}^{0}$} \\
\midrule
\ce{Pt6Ge4CO} & 2 & -1.947 & 3.837 & 2.545 & 1.292 & 0.337 \\
\ce{Pt6Ge4CO} & 3 & -1.732 & 3.961 & 1.730 & 2.231 & 0.563 \\
\ce{Pt6Ge4CO} & 1 & -1.657 & 4.869 & 2.569 & 2.300 & 0.472 \\
\ce{Pt6Ge4CO} & 6 & -1.842 & 3.952 & 1.752 & 2.200 & 0.557 \\
\ce{Pt6Ge4CO} & 5 & -1.215 & 3.889 & 0.770 & 3.119 & 0.802 \\
\ce{Pt6Ge4CO} & 4 & -1.207 & 3.870 & 0.748 & 3.122 & 0.807 \\
\midrule
\ce{Pt5Ge5CO} &  1/2/4/5 & -1.346 & 3.482 & 0.242 & 3.240 & 0.930 \\
\ce{Pt5Ge5CO} & 3 & -0.788 & 4.029 & 0.009 & 4.020 & 0.998 \\
\midrule
\ce{Pt4Ge6CO} & 2 & -1.506 & 4.609 & 0.370 & 4.239 & 0.920 \\
\ce{Pt4Ge6CO} & 4 & -1.240 & 3.015 & 0.004 & 3.010 & 0.999 \\
\ce{Pt4Ge6CO} & 1 & -1.150 & 2.992 & 0.002 & 2.989 & 0.999 \\
\ce{Pt4Ge6CO} & 3 & -1.046 & 4.710 & 0.447 & 4.263 & 0.905 \\
\bottomrule
\end{tabular}
\end{center}
\end{table}

\begin{table}[H]
\begin{center}
\caption{Local Hoppe ECoN descriptors for the Pt atom bound to \ce{CO} in
the optimized crossed-parent Pt--Ge carbonyls. The superscript \ce{CO}
distinguishes these final-state descriptors from the pre-adsorption values in
Table~\ref{tab:tfg_crossed_ptge_econ_descriptors}. Only metal neighbors are
included, with
$ECoN_{\rm total}^{\ce{CO}}=ECoN_{\rm Pt}^{\ce{CO}}+ECoN_{\rm Ge}^{\ce{CO}}$
and
$x_{\rm Ge}^{\ce{CO}}=ECoN_{\rm Ge}^{\ce{CO}}/ECoN_{\rm total}^{\ce{CO}}$.}
\label{tab:tfg_crossed_ptge_econ_co_descriptors}
\small
\begin{tabular}{lcrrrrr}
\toprule
{System} & {Site} & {$BE$ / eV} &
{$ECoN_{\rm total}^{\ce{CO}}$} & {$ECoN_{\rm Pt}^{\ce{CO}}$} &
{$ECoN_{\rm Ge}^{\ce{CO}}$} & {$x_{\rm Ge}^{\ce{CO}}$} \\
\midrule
\ce{Pt6Ge4CO}  & 2   & -2.246 & 3.004 & 1.774 & 1.230 & 0.409 \\
\ce{Pt6Ge4CO}  & 9   & -2.245 & 3.012 & 1.163 & 1.849 & 0.614 \\
\ce{Pt6Ge4CO}  & 1/4 & -2.235 & 4.300 & 2.425 & 1.876 & 0.436 \\
\ce{Pt6Ge4CO}  & 5/7 & -1.567 & 2.873 & 0.657 & 2.216 & 0.771 \\
\midrule
\ce{Pt5Ge5CO}  & 6/8  & -2.076 & 5.691 & 1.688 & 4.002 & 0.703 \\
\ce{Pt5Ge5CO}  & 7/10 & -1.661 & 2.952 & 0.937 & 2.015 & 0.683 \\
\ce{Pt5Ge5CO}  & 4    & -1.342 & 2.980 & 0.000 & 2.980 & 1.000 \\
\midrule
\ce{Pt4Ge6CO}  & 5/7 & -1.459 & 3.920 & 0.686 & 3.234 & 0.825 \\
\ce{Pt4Ge6CO}  & 2   & -1.305 & 3.221 & 0.134 & 3.086 & 0.958 \\
\ce{Pt4Ge6CO}  & 1   & -1.052 & 5.902 & 1.423 & 4.479 & 0.759 \\
\bottomrule
\end{tabular}
\end{center}
\end{table}

\begin{table}[H]
\begin{center}
\caption{Local Hoppe ECoN descriptors for the Pt atom bound to \ce{CO} in
the optimized GM-parent Pt--Ge carbonyls. Each composition is represented
using the spin state of its bare-cluster GM. \protect\blue{The GM is a triplet for
\ce{Pt6Ge4} and a singlet for \ce{Pt5Ge5} and \ce{Pt4Ge6}.} The superscript \ce{CO} indicates that these
descriptors are evaluated after adsorption and relaxation.}
\label{tab:tfg_gm_ptge_econ_descriptors}
%\scriptsize
\begin{tabular}{llrrrrr}
\toprule
{System} & {Site} & {$BE$ / eV} &
{$ECoN_{\rm total}^{\ce{CO}}$} & {$ECoN_{\rm Pt}^{\ce{CO}}$} &
{$ECoN_{\rm Ge}^{\ce{CO}}$} & {$x_{\rm Ge}^{\ce{CO}}$} \\
\midrule
\ce{Pt6Ge4CO} & 2  & -1.947 & 3.968 & 2.849 & 1.119 & 0.282 \\
\ce{Pt6Ge4CO} & 6  & -1.842 & 4.283 & 2.354 & 1.929 & 0.450 \\
\ce{Pt6Ge4CO} & 3  & -1.732 & 4.098 & 2.146 & 1.952 & 0.476 \\
\ce{Pt6Ge4CO} & 1  & -1.657 & 5.077 & 3.032 & 2.044 & 0.403 \\
\ce{Pt6Ge4CO} & 5  & -1.215 & 3.916 & 0.895 & 3.021 & 0.771 \\
\ce{Pt6Ge4CO} & 4  & -1.207 & 2.931 & 0.774 & 2.157 & 0.736 \\
\midrule
\ce{Pt5Ge5CO} & 1/2/4/5  & -1.346 & 4.523 & 1.547 & 2.976 & 0.658 \\
\ce{Pt5Ge5CO} & 3  & -0.788 & 3.962 & 0.005 & 3.957 & 0.999 \\
\midrule
\ce{Pt4Ge6CO} & 2  & -1.506 & 4.605 & 0.406 & 4.199 & 0.912 \\
\ce{Pt4Ge6CO} & 4  & -1.240 & 2.968 & 0.001 & 2.967 & 1.000 \\
\ce{Pt4Ge6CO} & 1  & -1.150 & 4.082 & 0.220 & 3.861 & 0.946 \\
\ce{Pt4Ge6CO} & 3  & -1.046 & 4.705 & 0.550 & 4.155 & 0.883 \\
\bottomrule
\end{tabular}
\end{center}
\end{table}
